\chapter{Endometrial Ablation}

\section*{Overview}
Endometrial ablation is a minimally invasive procedure that destroys the lining of the uterus (endometrium) to treat abnormal uterine bleeding in women who do not wish to preserve fertility. The exact approach, setting, and preparation depend on the indication, anatomy, and the patient's health.

\section*{Alternative Names}
Ablation of endometrium; endometrial resection; global endometrial ablation; NovaSure ablation

\section*{Principle}
Thermal energy (radiofrequency, microwave, balloon heat, or cryotherapy) is applied to the endometrial surface, destroying the lining to a controlled depth. Scarring of the cavity reduces or eliminates menstrual bleeding.

\section*{History}
Endometrial ablation was introduced with laser and resection techniques in the 1980s. Global ablation devices were approved in the late 1990s and early 2000s, simplifying the procedure.

\section*{Why It Is Done}
Performed for heavy menstrual bleeding (menorrhagia) that has failed medical therapy and is confirmed benign, in women who have completed childbearing.

\section*{Is This a Primary or Secondary Test or Procedure}
Primary alternative to hysterectomy for benign heavy menstrual bleeding.

\section*{How It Is Performed}
Performed as an outpatient procedure, often under local anesthesia or sedation. The cervix is dilated, the cavity is measured and inspected, and the ablation device is inserted and activated for several minutes, delivering energy that destroys the endometrium.

\section*{Preparation}
Preprocedure endometrial sampling to exclude hyperplasia or cancer, exclusion of pregnancy, and confirmation that childbearing is complete. Endometrial thinning agents may be used.

\section*{Contraindications and Precautions}
Pregnancy, uterine infection, known or suspected uterine cancer, recent pregnancy, and desire for future fertility are contraindications. Precaution: uterine anomalies or a very thin myometrium increase perforation risk.

\section*{Risks}
Uterine perforation, bleeding, infection, thermal injury to adjacent organs, hematometra, and postablation syndrome with cyclic pain.

\section*{Aftercare and Recovery}
Mild cramping and watery discharge for a few weeks. Birth control is still recommended if there is a risk of pregnancy because conception can be dangerous.

\section*{Outcome and Follow-up}
The procedure eliminates or markedly reduces menstrual bleeding in most women; a minority become amenorrheic. Bleeding may take several months to settle.

\section*{Expected Results}
A significant reduction in menstrual flow with no complications, adequate for most women to avoid hysterectomy.

\section*{Follow-up}
Review at 3-6 months for bleeding outcomes; persistent bleeding requires further evaluation for endometrial pathology.

\section*{When to Seek Medical Care}
Follow the procedure team's specific instructions. Seek urgent medical assessment for severe or worsening pain, uncontrolled bleeding, fever or signs of infection, trouble breathing, chest pain, fainting, new weakness or confusion, inability to urinate, or any rapidly worsening symptom. The appropriate warning signs depend on the procedure and the patient's medical condition.

\section*{References}
\begin{enumerate}
  \item MedlinePlus. Endometrial ablation. U.S. National Library of Medicine; 2025.
  \item Current Medical Diagnosis and Treatment. McGraw-Hill; 2025.
\end{enumerate}

\clearpage
